% ৩.৩ সংখ্যা পদ্ধতির রূপান্তর
\section{সংখ্যা রূপান্তর (Conversion)}

\subsection{Decimal থেকে Binary}
\begin{itemize}
\item ভাগ ও remainder পদ্ধতি
\item ভগ্নাংশ রূপান্তরের সারমর্ম
\end{itemize}

\subsection{Binary থেকে Decimal}
\begin{itemize}
\item positional value এবং 2-এর ঘাত ব্যবহার
\end{itemize}

\subsection{Decimal ও Octal/Hexadecimal এবং Binary ও Octal/Hex মধ্যে রূপান্তর}
\begin{itemize}
\item Decimal থেকে Octal/Hex: পুনরাবৃত্ত ভাগ
\item Binary থেকে Octal: 3-bit grouping
\item Binary থেকে Hex: 4-bit grouping
\item Octal/Hex এর মধ্যে রূপান্তর: Binary intermediary
\end{itemize}

\begin{learningbox}
এই টপিক শেষে শিক্ষার্থী পারবে:
\begin{itemize}
\item decimal থেকে binary/octal/hex রূপান্তর করতে।
\item binary থেকে decimal, octal ও hexadecimal রূপান্তর করতে।
\item fraction রূপান্তরে multiplication method-এর ধারণা ব্যবহার করতে।
\end{itemize}
\end{learningbox}

\begin{keypointbox}{রূপান্তরের shortcut}
\begin{itemize}
\item Decimal থেকে অন্য base: base দিয়ে বারবার ভাগ, remainder নিচ থেকে ওপর।
\item অন্য base থেকে Decimal: প্রতিটি digit $\times$ base-এর ঘাত যোগ।
\item Binary থেকে Octal: ডান দিক থেকে 3-bit group।
\item Binary থেকে Hex: ডান দিক থেকে 4-bit group।
\end{itemize}
\end{keypointbox}

\begin{examplebox}{Decimal থেকে Binary}
$(25)_{10}$ রূপান্তর: $25\div2=12$ remainder 1, $12\div2=6$ remainder 0, $6\div2=3$ remainder 0, $3\div2=1$ remainder 1, $1\div2=0$ remainder 1। তাই $(25)_{10}=(11001)_2$।
\end{examplebox}

\begin{examplebox}{Binary থেকে Decimal}
$(10110)_2=1\times2^4+0\times2^3+1\times2^2+1\times2^1+0\times2^0=16+4+2=22$।
\end{examplebox}

\begin{center}
\fbox{\parbox{0.9\textwidth}{
\centering
\textbf{Conversion flowchart placeholder}\\[4pt]
Decimal $\leftrightarrow$ Binary $\leftrightarrow$ Octal/Hex\\
Octal ও Hex-এর মধ্যে সরাসরি না করে binary intermediary ব্যবহার দেখাতে হবে।
}}
\end{center}

\begin{practicebox}
\textbf{অনুশীলন}
\begin{enumerate}
\item $(43)_{10}$ কে binary-তে রূপান্তর কর।
\item $(101101)_2$ কে decimal-এ রূপান্তর কর।
\item $(11011101)_2$ কে hexadecimal-এ রূপান্তর কর।
\end{enumerate}
\end{practicebox}
